\documentclass{article}
\usepackage[margin=1in]{geometry}
\usepackage[T1]{fontenc}
\usepackage{amsmath,amssymb}
\newtheorem{proposition}{Proposition}
\newcommand{\norm}[1]{\left\|#1\right\|}
\title{Conditional closed-loop gradient and finite-difference error bounds}
\author{EAGM 2.0 research notes; elementary conditional results}
\date{November 2026}
\begin{document}\maketitle
\section{A forward-error non-implication}
Set $F(x)=0$ and $\widehat F_\epsilon(x)=\epsilon\sin(x/\epsilon^2)$. Uniform forward error is at most $\epsilon$, but derivative error at zero is $1/\epsilon$. Both maps are smooth. Therefore forward error alone, including a small empirical mean squared error, cannot control Jacobian error without additional uniform regularity assumptions. This is a counterexample, not a diagnosis that the fitted neural networks realize this oscillation.

\section{Coordinates and explicit assumptions}
Let $z$ include physical, execution, previous-action and reference information sufficient to be a complete state. For fixed discrete execution configuration, define
$T_\theta(z)=F(z,\pi_\theta(z))$ and $\widehat T_\theta(z)=\widehat F(z,\pi_\theta(z))$.
Both use the same coordinates and policy. Any latent model application requires an independently justified differentiable representation/lifting; the deployed GRU is not automatically such a representation. Constants below are uniform on a convex neighborhood containing both rollouts and the connecting segments, for the parameter positions considered. Assume $F,\widehat F,\pi$ differentiable, with
\begin{align*}
\norm{F_z},\norm{\widehat F_z}&\le L_z,&
\norm{F_a},\norm{\widehat F_a}&\le L_a,&
\norm{\pi_z}&\le P_z,&\norm{\pi_\theta}&\le P_\theta,\\
\norm{\widehat F-F}&\le\epsilon_F,&
\norm{\widehat F_z-F_z}&\le\epsilon_z,&
\norm{\widehat F_a-F_a}&\le\epsilon_a.
\end{align*}
Use joint distance $\norm{z-z'}+\norm{a-a'}$; suppose $F_z,F_a$ are Lipschitz with constants $K_z,K_a$, and $\pi_z,\pi_\theta$ are Lipschitz in $z$ with constants $Q_z,Q_\theta$. Then valid constants for $T_z,T_\theta$ are
\begin{align*}
L&=L_z+L_aP_z,& B&=L_aP_\theta,\\
K_A&=(K_z+K_aP_z)(1+P_z)+L_aQ_z,&
K_B&=K_a(1+P_z)P_\theta+L_aQ_\theta,\\
\epsilon_A&=\epsilon_z+\epsilon_aP_z,&
\epsilon_B&=\epsilon_aP_\theta.
\end{align*}
\textit{Justification.} The chain rule gives $T_z=F_z+F_a\pi_z$ and $T_\theta=F_a\pi_\theta$. Bound each product difference by adding and subtracting one common-factor product. The policy changes the joint input distance by at most $1+P_z$. At the same state the policy factors cancel, giving $\epsilon_A,\epsilon_B$. Bounds for both model and true derivatives give the common $L$.

\section{State-estimation and sensitivity errors}
Let $e_t=\norm{\widehat z_t-z_t}$, $S_t=\partial z_t/\partial\theta$, $\widehat S_t=\partial\widehat z_t/\partial\theta$, $s_t=\norm{S_t}$ and $E_t=\norm{\widehat S_t-S_t}$. Initial history estimation may have both $e_0>0$ and $E_0>0$.
\begin{proposition}
Under the assumptions above,
\begin{align}
e_{t+1}&\le Le_t+\epsilon_F, &s_{t+1}&\le Ls_t+B,\\
E_{t+1}&\le LE_t+(\epsilon_A+K_Ae_t)s_t+\epsilon_B+K_Be_t.\label{eq:rec}
\end{align}
Consequently $e_t\le L^te_0+\epsilon_F\sum_{i=0}^{t-1}L^i$, $s_t\le L^ts_0+B\sum_{i=0}^{t-1}L^i$, and
\[
E_t\le L^tE_0+\sum_{k=0}^{t-1}L^{t-1-k}\{(\epsilon_A+K_Ae_k)s_k+\epsilon_B+K_Be_k\}.
\]
\end{proposition}
\noindent\textit{Proof.} Add and subtract $F(\widehat z,\pi(\widehat z))$ for the state bound. The chain rule gives $S_{t+1}=A_tS_t+B_t$ and its model counterpart. Subtract them as
$\widehat A_t(\widehat S_t-S_t)+(\widehat A_t-A_t)S_t+\widehat B_t-B_t$.
At their respective states the derivative discrepancies are bounded by $\epsilon_A+K_Ae_t$ and $\epsilon_B+K_Be_t$. Apply triangle and operator-norm inequalities, then induction. $\square$

If the saved physical start and observed history are fixed while optimizing theta, both initial sensitivities are zero: $s_0=E_0=0$. This does not set $e_0=0$ for hidden-state estimation. If history is itself generated by theta, its derivative cannot be discarded. The implementation's frozen starts analyze a conditional rollout objective, not derivatives of its data-collection distribution.

\section{Reward and policy regularity}
Write the per-step scalar return as $R(z,a)$ and its modeled counterpart $\widehat R(z,a)$; these may absorb next-state rewards through $F$ and previous action through $z$. Assume $\norm{\widehat R_z}\le R_z^*$, $\norm{\widehat R_a}\le R_a^*$ and the same true bounds, with same-input derivative errors $\epsilon_{Rz},\epsilon_{Ra}$. Suppose true $R_z,R_a$ have joint Lipschitz constants $K_{Rz},K_{Ra}$. For closed-loop $r_\theta(z)=R(z,\pi_\theta(z))$, define
\begin{align*}
C&=R_z^*+R_a^*P_z,&
\epsilon_C&=\epsilon_{Rz}+\epsilon_{Ra}P_z,&
\epsilon_D&=\epsilon_{Ra}P_\theta,\\
K_C&=(K_{Rz}+K_{Ra}P_z)(1+P_z)+R_a^*Q_z,&
K_D&=K_{Ra}(1+P_z)P_\theta+R_a^*Q_\theta.
\end{align*}
For $J=\sum_{t=0}^{H-1}\gamma^t r_\theta(z_t)$ and $0\le\gamma\le1$,
\begin{equation}
\norm{\nabla\widehat J-\nabla J}\le\sum_{t=0}^{H-1}\gamma^t\left[CE_t+(\epsilon_C+K_Ce_t)s_t+\epsilon_D+K_De_t\right].
\end{equation}
\textit{Proof.} The contribution is $C_tS_t+D_t$ with $C_t=\partial_zr_\theta$, $D_t=\partial_\theta r_\theta$. Subtract modeled and true contributions, adding and subtracting $\widehat C_tS_t$. The reward product bounds follow identically to the transition ones; sum the triangle inequalities. $\square$

For a learned visual reward, errors $\epsilon_{Rz},\epsilon_{Ra}$ are not zero by definition. Even for an analytic physical reward, composing it with an imperfect transition creates reward derivative discrepancy. Changing the start distribution adds a separate gradient-distribution error. The bound can grow exponentially when $L>1$; no measured or certified constants are supplied, so it is not a numerical performance guarantee.

\section{Finite differences matching implemented stencils}
For coordinate restriction $f(s)=J(\theta+se_j)$ assume $f\in C^3$ on the full feasible stencil interval, $|f'''|\le M_3$, and each reported objective has error at most $\eta$. These are deterministic local conditions, not inferred from scale agreement.
\begin{proposition}
For symmetric central differences at $\pm h$, error is at most $M_3h^2/6+\eta/h$. For unequal offsets $-a,0,b$ with $a,b>0$, weights
\[
w_-=-\frac{b}{a(a+b)},\quad w_0=\frac{b-a}{ab},\quad w_+=\frac{a}{b(a+b)}
\]
give error at most $M_3ab/6+\eta(|w_-|+|w_0|+|w_+|)$. At the exact lower boundary the stencil $(-3f(0)+4f(h)-f(2h))/(2h)$ has error at most $M_3h^2/3+4\eta/h$; reflection gives the upper-boundary result.
\end{proposition}
\noindent\textit{Proof.} Taylor expansion cancels constant and quadratic terms, while the weighted linear term equals $f'(0)$. For the unequal case the third-order remainder bound is
$M_3[ba^2+ab^2]/[6(a+b)]=M_3ab/6$. The symmetric case follows by $a=b=h$. For the forward case use the integral Taylor remainder. Its kernel is $[4(h-s)^2-(2h-s)^2]/(4h)$ for $0\le s\le h$ and $-(2h-s)^2/(4h)$ for $h\le s\le2h$. Both are nonpositive and their total integral is $-h^2/3$. Noise contributes the sum of absolute weights times $\eta$. $\square$

The implementation reduces $h$ if needed to keep both one-sided points feasible, and includes cached center queries in accounting. Clipping a symmetric stencil without correcting its weights does not satisfy this result. At a contact or clipping kink, $C^3$ fails and multi-scale agreement cannot repair the missing hypothesis. With stochastic rollouts an almost-sure bounded-error or separate probabilistic noise model is needed; these inequalities alone give no confidence interval.

\section{Actual policy update, not raw-gradient rhetoric}
If the evaluation objective has $L_J$-Lipschitz gradient on the actual update segment and $\norm{\widehat g-\nabla J}\le\delta$, then for any feasible actual step $d$,
\[
J(\theta+d)-J(\theta)\ge\widehat g^Td-\delta\norm d-\frac{L_J}{2}\norm d^2.
\]
Integrate the gradient along the segment and apply Cauchy--Schwarz to prove this inequality. Adam and projection require their actual resulting $d$; a high cosine alone neither supplies $\delta$ nor a safe step size. No convergence, improvement, contact differentiability, latent sufficiency or novelty claim follows without these additional assumptions.
\end{document}
